\section{Conclusion}

We addressed the problem of maintaining factual truth in long-context
conversational memory. The task is not merely to retrieve textually
similar past utterances, but to ensure that the information presented
to the model over time remains consistent with the evolving ground
truth, even as user statements are corrected or superseded.

Our solution introduces a structured memory mechanism that explicitly
models the truth status of information. Each memory unit is a claim
schema annotated with its validity and linked to others via
\textit{contradict} and \textit{supersede} edges. This graph enables
truth-consistent retrieval, where only currently valid claims are
surfaced, and strategic abstention when the memory lacks sufficient
information to answer reliably. This approach directly enforces a
coherence constraint over the memory's contents.

Empirically, TTMG achieves paired set-inclusion gains on two
truth-maintenance-sensitive capability slices -- abstention
(+\DeltaLMEABS{}\,pp; paired 4--0) and single-session-user
(+\DeltaLMESSU{}\,pp; paired 2--0) -- at \TokensOverhead\% of A-Mem's
token cost. On the unselected bulk of the benchmark, TTMG and Flat
RAG are statistically indistinguishable (overall $p{=}$\McnemarOverall);
composing the two via inference-time routers (TTMG-Abstain
\AbstainRouterAcc\%, TTMG+R-LLM \LLMRouterAcc\%) matches Flat RAG's
overall accuracy while preserving TTMG's abstention and supersede
behaviours. The oracle upper bound of \OracleLMEAcc\% documents
that claim-level and raw-text retrieval carry substantially
complementary information; harvesting
that complementarity beyond the feature router is a direct target for
future work. Explicitly reasoning about truth maintenance is therefore
a useful primitive to add to long-running conversational agents,
particularly where abstention and user-fact fidelity matter more than
bulk retrieval accuracy.
